% GENERATED FILE. Edit canonical lessons, then run npm run generate:books.
% canonical-source-hash: fnv1a64:4e6efff6274d2bc2
% canonical-lessons: TE-C286-vayida-veyu, TE-C286-venakadu, TE-C286-tappincuko, TE-C286-varcu, TE-C286-anuko, TE-R286-first-pass-words-from-chapters-278-281, TE-R286-second-pass-words-from-chapters-282-286

\chapter{To postpone, To hesitate, To escape, To drain, To lean}
\label{ch:te-to-postpone-to-hesitate-to-escape-to-drain-to-lean}
\hlchaptermodality{286}

\begin{chapteropening}
\textbf{By the end of this chapter:} \emph{I can say to postpone, to hesitate, to escape, to drain and to lean.}

Complete the last lesson of chapter 286: I can say to postpone, to hesitate, to escape, to drain and to lean.
\end{chapteropening}

\section[\texorpdfstring{\te{వాయిదా} \te{వేయు} (vāyidā vēyu)}{వాయిదా వేయు (vāyidā vēyu)}]{\te{వాయిదా} \te{వేయు} (vāyidā vēyu) — to postpone}

\label{lesson:TE-C286-vayida-veyu}



\begin{quote}
Before the new one: say the Telugu for to chase, then the Telugu for to melt.
\end{quote}

\subsection*{You'll want to know: \te{వాయిదా} \te{వేయు}}
\textbf{\te{వాయిదా} \te{వేయు}} — \emph{vāyidā vēyu} — \textquotedblleft{}to postpone\textquotedblright{}.

\begin{grammarlens}[title={What you've built}]
One more word for this chapter.
\end{grammarlens}

\subsection*{Your turn}
\begin{itemize}
\raggedright
  \item \emph{Say it:} \emph{vāyidā vēyu}
  \item \emph{Say it:} \emph{vāyidā vēyu}, once more
  \item \emph{Say it:} say \emph{vāyidā vēyu}
  \item \emph{Recall:} say the Telugu for to chase, then the Telugu for to melt, then say \emph{vāyidā vēyu} again
\end{itemize}

\begin{tcolorbox}[breakable,colback=teal!4,colframe=teal!35!black,title={Before you move on}]
What does \te{వాయిదా} \te{వేయు} mean? (\textquotedblleft{}To postpone\textquotedblright{}.) Say it once more.
\end{tcolorbox}
\section[\texorpdfstring{\te{వెనకాడు} (venakāḍu)}{వెనకాడు (venakāḍu)}]{\te{వెనకాడు} (venakāḍu) — to hesitate}

\label{lesson:TE-C286-venakadu}



\begin{quote}
Before the new one: say the Telugu for to melt, then the Telugu for to postpone.
\end{quote}

\subsection*{You'll want to know: \te{వెనకాడు}}
\textbf{\te{వెనకాడు}} — \emph{venakāḍu} — \textquotedblleft{}to hesitate\textquotedblright{}.

\begin{grammarlens}[title={What you've built}]
One more word for this chapter.
\end{grammarlens}

\subsection*{Your turn}
\begin{itemize}
\raggedright
  \item \emph{Say it:} \emph{venakāḍu}
  \item \emph{Say it:} \emph{venakāḍu}, once more
  \item \emph{Say it:} say \emph{venakāḍu}
  \item \emph{Recall:} say the Telugu for to melt, then the Telugu for to postpone, then say \emph{venakāḍu} again
\end{itemize}

\begin{tcolorbox}[breakable,colback=teal!4,colframe=teal!35!black,title={Before you move on}]
What does \te{వెనకాడు} mean? (\textquotedblleft{}To hesitate\textquotedblright{}.) Say it once more.
\end{tcolorbox}
\section[\texorpdfstring{\te{తప్పించుకో} (tappiñcukō)}{తప్పించుకో (tappiñcukō)}]{\te{తప్పించుకో} (tappiñcukō) — to escape}

\label{lesson:TE-C286-tappincuko}



\begin{quote}
Before the new one: say the Telugu for to postpone, then the Telugu for to hesitate.
\end{quote}

\subsection*{You'll want to know: \te{తప్పించుకో}}
\textbf{\te{తప్పించుకో}} — \emph{tappiñcukō} — \textquotedblleft{}to escape\textquotedblright{}.

\begin{grammarlens}[title={What you've built}]
One more word for this chapter.
\end{grammarlens}

\subsection*{Your turn}
\begin{itemize}
\raggedright
  \item \emph{Say it:} \emph{tappiñcukō}
  \item \emph{Say it:} \emph{tappiñcukō}, once more
  \item \emph{Say it:} say \emph{tappiñcukō}
  \item \emph{Recall:} say the Telugu for to postpone, then the Telugu for to hesitate, then say \emph{tappiñcukō} again
\end{itemize}

\begin{tcolorbox}[breakable,colback=teal!4,colframe=teal!35!black,title={Before you move on}]
What does \te{తప్పించుకో} mean? (\textquotedblleft{}To escape\textquotedblright{}.) Say it once more.
\end{tcolorbox}
\section[\texorpdfstring{\te{వార్చు} (vārcu)}{వార్చు (vārcu)}]{\te{వార్చు} (vārcu) — to drain (cooked rice)}

\label{lesson:TE-C286-varcu}



\begin{quote}
Before the new one: say the Telugu for to hesitate, then the Telugu for to escape.
\end{quote}

\subsection*{You'll want to know: \te{వార్చు}}
\textbf{\te{వార్చు}} — \emph{vārcu} — \textquotedblleft{}to drain (cooked rice)\textquotedblright{}.

\begin{grammarlens}[title={What you've built}]
One more word for this chapter.
\end{grammarlens}

\subsection*{Your turn}
\begin{itemize}
\raggedright
  \item \emph{Say it:} \emph{vārcu}
  \item \emph{Say it:} \emph{vārcu}, once more
  \item \emph{Say it:} say \emph{vārcu}
  \item \emph{Recall:} say the Telugu for to hesitate, then the Telugu for to escape, then say \emph{vārcu} again
\end{itemize}

\begin{tcolorbox}[breakable,colback=teal!4,colframe=teal!35!black,title={Before you move on}]
What does \te{వార్చు} mean? (\textquotedblleft{}To drain (cooked rice)\textquotedblright{}.) Say it once more.
\end{tcolorbox}
\section[\texorpdfstring{\te{ఆనుకో} (ānukō)}{ఆనుకో (ānukō)}]{\te{ఆనుకో} (ānukō) — to lean (against)}

\label{lesson:TE-C286-anuko}



\begin{quote}
Before the new one: say the Telugu for to escape, then the Telugu for to drain.
\end{quote}

\subsection*{You'll want to know: \te{ఆనుకో}}
\textbf{\te{ఆనుకో}} — \emph{ānukō} — \textquotedblleft{}to lean (against)\textquotedblright{}.

\begin{grammarlens}[title={What you've built}]
One more word for this chapter.
\end{grammarlens}

\subsection*{Your turn}
\begin{itemize}
\raggedright
  \item \emph{Say it:} \emph{ānukō}
  \item \emph{Say it:} \emph{ānukō}, once more
  \item \emph{Say it:} say \emph{ānukō}
  \item \emph{Recall:} say the Telugu for to escape, then the Telugu for to drain, then say \emph{ānukō} again
\end{itemize}

\begin{tcolorbox}[breakable,colback=teal!4,colframe=teal!35!black,title={Before you move on}]
What does \te{ఆనుకో} mean? (\textquotedblleft{}To lean (against)\textquotedblright{}.) Say it once more.
\end{tcolorbox}
\section[(dialogue)]{Words from chapters 278-281 — a first pass}

\label{lesson:TE-R286-first-pass-words-from-chapters-278-281}



\begin{quote}
Say the words for to drain and to lean.
\end{quote}

\subsection*{Your turn}
Say these aloud:

\begin{itemize}
\raggedright
  \item to hold, to bathe, to prepare, to cross, to drive
  \item to recover, to swell, to feel dizzy, to touch, to peel
  \item to rub, to receive, to tell, to admit, to deposit
  \item to cancel, to submit, to press, to get lost, to love
\end{itemize}

\begin{tcolorbox}[breakable,colback=teal!4,colframe=teal!35!black,title={Before you move on}]
To hold? (\textbf{\te{పట్టుకో}}, \emph{paṭṭukō}.) To bathe? (\textbf{\te{స్నానం} \te{చేయు}}, \emph{snānaṁ cēyu}.) To prepare? (\textbf{\te{తయారు} \te{చేయు}}, \emph{tayāru cēyu}.) To cross? (\textbf{\te{దాటు}}, \emph{dāṭu}.) To drive? (\textbf{\te{నడుపు}}, \emph{naḍupu}.) To recover? (\textbf{\te{కోలుకో}}, \emph{kōlukō}.) To swell? (\textbf{\te{ఉబ్బు}}, \emph{ubbu}.) To feel dizzy? (\textbf{\te{కళ్ళు} \te{తిరుగు}}, \emph{kaḷḷu tirugu}.) To touch? (\textbf{\te{తాకు}}, \emph{tāku}.) To peel? (\textbf{\te{ఒలుచు}}, \emph{olucu}.) To rub? (\textbf{\te{రుద్దు}}, \emph{ruddu}.) To receive? (\textbf{\te{అందుకో}}, \emph{andukō}.) To tell? (\textbf{\te{చెప్పు}}, \emph{ceppu}.) To admit? (\textbf{\te{చేర్చు}}, \emph{cērcu}.) To deposit? (\textbf{\te{జమ} \te{చేయు}}, \emph{jama cēyu}.) To cancel? (\textbf{\te{రద్దు} \te{చేయు}}, \emph{raddu cēyu}.) To submit? (\textbf{\te{సమర్పించు}}, \emph{samarpiñcu}.) To press? (\textbf{\te{నొక్కు}}, \emph{nokku}.) To get lost? (\textbf{\te{తప్పిపోవు}}, \emph{tappipōvu}.) To love? (\textbf{\te{ప్రేమించు}}, \emph{prēmiñcu}.)
\end{tcolorbox}
\section[(dialogue)]{Words from chapters 282-286 — a second pass}

\label{lesson:TE-R286-second-pass-words-from-chapters-282-286}



\begin{quote}
Say the words for to graze and to plant.
\end{quote}

\subsection*{Your turn}
Say these aloud:

\begin{itemize}
\raggedright
  \item to graze, to plant, to decorate, to share out, to visit
  \item to charge, to remove, to recognize, to oppose, to argue
  \item to let go, to dig, to happen, to sweep, to knead
  \item to threaten, to warn, to take note of, to chase, to melt
  \item to postpone, to hesitate, to escape, to drain, to lean
\end{itemize}

\begin{tcolorbox}[breakable,colback=teal!4,colframe=teal!35!black,title={Before you move on}]
To graze? (\textbf{\te{మేయు}}, \emph{mēyu}.) To plant? (\textbf{\te{నాటు}}, \emph{nāṭu}.) To decorate? (\textbf{\te{అలంకరించు}}, \emph{alaṅkariñcu}.) To share out? (\textbf{\te{పంచు}}, \emph{pañcu}.) To visit? (\textbf{\te{సందర్శించు}}, \emph{sandarśiñcu}.) To charge? (\textbf{\te{ఛార్జ్} \te{చేయు}}, \emph{chārj cēyu}.) To remove? (\textbf{\te{తొలగించు}}, \emph{tolagiñcu}.) To recognize? (\textbf{\te{గుర్తుపట్టు}}, \emph{gurtupaṭṭu}.) To oppose? (\textbf{\te{వ్యతిరేకించు}}, \emph{vyatirēkiñcu}.) To argue? (\textbf{\te{వాదించు}}, \emph{vādiñcu}.) To let go? (\textbf{\te{వదులు}}, \emph{vadulu}.) To dig? (\textbf{\te{తవ్వు}}, \emph{tavvu}.) To happen? (\textbf{\te{జరుగు}}, \emph{jarugu}.) To sweep? (\textbf{\te{చిమ్ము}}, \emph{cimmu}.) To knead? (\textbf{\te{పిసుకు}}, \emph{pisuku}.) To threaten? (\textbf{\te{బెదిరించు}}, \emph{bediriñcu}.) To warn? (\textbf{\te{హెచ్చరించు}}, \emph{heccariñcu}.) To take note of? (\textbf{\te{గమనించు}}, \emph{gamaniñcu}.) To chase? (\textbf{\te{తరుము}}, \emph{tarumu}.) To melt? (\textbf{\te{కరుగు}}, \emph{karugu}.) To postpone? (\textbf{\te{వాయిదా} \te{వేయు}}, \emph{vāyidā vēyu}.) To hesitate? (\textbf{\te{వెనకాడు}}, \emph{venakāḍu}.) To escape? (\textbf{\te{తప్పించుకో}}, \emph{tappiñcukō}.) To drain? (\textbf{\te{వార్చు}}, \emph{vārcu}.) To lean? (\textbf{\te{ఆనుకో}}, \emph{ānukō}.)
\end{tcolorbox}
